\documentclass[a4paper,11pt]{article}
\pagestyle{empty}
\title{Notes sur le Lab 1}
\author{Antoine POIRIER ING5 DEF}
\begin{document}

\maketitle

\section*{O\`u peuvent s'infiltrer les instructions douteuses ?}

Ces instructions peuvent s'imiscer dans le système IA de deux fa\c{c}ons principales : soit de manière directe gr\^ace à des prompts utilisateurs néfastes (direct prompt injection), soit indirectement par le biais des sources de données qui ne proviennent pas directement de l'utilisateur mais plut\^ot des plugins et/ou datasets utilis\'es par le LLM (indirection prompt injection).

\section*{Quelles mesures ont chang\'{e} le r\'{e}sultat ?}

La suppression d'expressions imp\'eratives et le r\'etrogradage des privil\`eges ont permit de r\'eduire les risques li\'es notamment \`a l'agence excessive du mod\`ele et aux injections.

\section*{Les risques identifi\'es}

\begin{tabular}{|c|c|c|}
	\hline
	Risque & OWASP & MITRE ALIAS \\
	Tentative injection pour obtenir les acc\`es admin. & LLM01 & AML.T0051 \\
	Affichage des authentifiants admin & LLM02 & AML.T0098  \\
	Traitement des sorties non-s\'ecuris\'e & LLM02 & AML.T0135 \\
	Privil\`eges excessives confi\'es au mod\`ele & LLM08 & AML.T0053 \\
	Tentative direct de jailbreak & LLM01 & AML.T0054 \\
	\hline
\end{tabular}

\end{document}

